La evolución de los sistemas de monitoreo ambiental ha transitado desde dispositivos analógicos aislados hacia ecosistemas complejos de \textit{Wireless Sensor Networks} (WSN) integrados en arquitecturas de Smart Buildings e Industria 4.0. El estado actual de la técnica se define por la capacidad de procesar datos en la periferia (\textit{Edge Computing}) y la interoperabilidad mediante protocolos ligeros.

\subsection{Sistemas SCADA y Monitoreo Industrial}
Históricamente, el monitoreo industrial ha dependido de sistemas de Control Supervisorio y Adquisición de Datos (SCADA). Sin embargo, los sistemas SCADA tradicionales suelen ser rígidos y centralizados. La tendencia moderna, impulsada por el IoT Industrial (IIoT), propone una transición hacia sistemas distribuidos donde cada nodo posee capacidad de cómputo independiente. Según \citet{valvano2014volume2}, el diseño de interfaces de tiempo real es el núcleo que permite a estos sistemas responder a eventos asíncronos con determinismo, superando la latencia de los sistemas de sondeo (\textit{polling}) convencionales.

\subsection{Redes de Sensores Inalámbricos (WSN) y MQTT}
El despliegue de redes \textit{NodeAlert IoT } se fundamenta en el concepto de WSN, donde la topología de red permite una cobertura extensa sin dependencia de cableado estructurado. El estándar de comunicación \textit{de facto} para estas redes es el protocolo MQTT (\textit{Message Queuing Telemetry Transport}). La literatura técnica destaca que la eficiencia de MQTT radica en su modelo de publicación/suscripción, ideal para dispositivos con recursos limitados de memoria y energía \citep{mazidi2014}. Esto permite la creación de sistemas de mensajería asíncrona que pueden escalar desde un entorno doméstico hasta una infraestructura de \textit{Smart Building}.

\subsection{Edge Computing y Sistemas Distribuidos Embebidos}
Una limitación crítica en el estado del arte previo era la dependencia de la nube para la toma de decisiones. El \textit{Edge Computing} resuelve esto procesando la lógica de control localmente en el nodo central o \textit{gateway}. \citet{valvano2014volume1} sostiene que la robustez de un sistema distribuido embebido reside en su capacidad de mantener la operatividad incluso ante la pérdida de conectividad externa. La integración de sistemas operativos de tiempo real (RTOS) en nodos periféricos como el ESP32 garantiza que la captura de señales críticas, como las descritas en la teoría de dispositivos electrónicos de \citet{boylestad2009}, no se vea interrumpida por las tareas de comunicación de red.

En este contexto, \textit{NodeAlert IoT} no solo actúa como un detector, sino como un sistema embebido distribuido que integra supervisión remota, alertas automáticas y lógica de control local, cerrando la brecha entre los prototipos académicos y las soluciones industriales escalables.